We begin by formalizing a model of a ReLU network, and defining various concepts related to polyhedra, PWA functions, and dynamical systems.
We then state the main problems we seek to address in this paper.
\subsection{ReLU Networks} \label{sec:relu_networks}
An $L$-layer feedforward neural network implements a function $\by = F_{\bTheta}(\bx)$, with the map $F_{\bTheta}: \R^{n} \rightarrow \R^{m}$ defined recursively by the iteration
\begin{align}
\label{Eq:NNBasic}
    \bz_{i} = \sigma_i(\bW_i \bz_{i-1} + \bb_i), \quad i = 1, \ldots, L
\end{align}
where $\bz_0 = \bx$ is the input, $\by = \bz_L$ is the output, and $\bz_i$ is the hidden layer output of the network at layer $i$.  The function $\sigma_i(\cdot)$ is the activation function at layer $i$, and $\bW_i\in\R^{l_i \times l_{i-1}}$ and $\bb_i\in\R^{l_i}$ are the weights and biases, respectively, where $l_i$ denotes the number of neurons in layer $i$. We assume $\sigma_{L}(\cdot)$ is an identity map and all hidden layer activations are ReLU,
\begin{align}
\label{Eq:ReLU}
    \sigma_i(\bz^\text{pre}_{i}) = [\max(0,z^\text{pre}_{i1}) \, \cdots \, \max(0,z^\text{pre}_{i l_i})]^\top,
\end{align}
where $\bz^\text{pre}_i = [z^\text{pre}_{i1} \, \cdots \, z^\text{pre}_{i l_i}]^\top = \bW_i\bz_{i-1} + \bb_i$ is the pre-activation hidden state at layer $i$. We can rewrite the iteration in (\ref{Eq:NNBasic}) in a homogeneous form
\begin{align}
\label{Eq:NNHomogeneous}
        \begin{bmatrix}
            \bz_i \\ 
            1
        \end{bmatrix}
            &= \sigma_i\Big(\bTheta_i \begin{bmatrix}
            \bz_{i-1} \\ 
            1
        \end{bmatrix}\Big),
\end{align}
where
\begin{subequations}
\begin{align}
\bTheta_i &= 
        \begin{bmatrix}
            \bW_i & \bb_i \\
            \mathbf{0}^\top & 1
        \end{bmatrix} \quad \text{for} \quad i = 1, \ldots, L-1, \label{Eq:ThetaDef} \\
\bTheta_L &= 
    \begin{bmatrix}
        \bW_L & \bb_L
    \end{bmatrix}. \label{Eq:ThetaLDef}
\end{align}
\end{subequations}
We define $\bTheta = (\bTheta_1, \ldots, \bTheta_L)$ as the tuple containing all the parameters that define the function $F_{\bTheta}(\cdot)$. 

Given an input $\bx$ to a ReLU network, every hidden neuron has an associated binary \textit{neuron activation} of zero or one, corresponding to whether the preactivation value is nonpositive or positive, respectively.\footnote{This choice is arbitrary. Some works use the opposite convention.} Specifically, the activation for neuron $j$ in layer $i$ is given by
\begin{align}
\label{Eq:AP}
    \lambda_{(ij)}(\bx) &= \begin{cases}
            1 \quad \text{for} \quad z^\text{pre}_{ij}(\bx) > 0 \\    
            0 \quad \text{for} \quad z^\text{pre}_{ij}(\bx) \le 0. 
                        \end{cases}
\end{align}
We define the activation pattern at layer $i$ as the binary vector $\blambda_{(i)}(\bx) = [\lambda_{(i1)}(\bx) \, \cdots \, \lambda_{(i l_i)}(\bx)]^\top$ and the activation pattern of the whole network as the tuple $\lambda(\bx) = (\blambda_{(1)}(\bx), \ldots, \blambda_{(L-1)}(\bx))$.  For a network with $N = \sum_{i = 1}^{L-1}l_i$ neurons there are $2^N$ possible combinations of neuron activations, however not all are realizable by the network due to inter-dependencies between neurons from layer to layer, as we will see later in the paper. Empirically, the number of activation patterns achievable by the network is better approximated by $N^{n}$ for input space with dimension $n$  \cite{hanin2019deep}. 

We can write equivalent expressions for the neural network function in terms of the activation pattern.  Define the diagonal activation matrix for layer $i$ as $\bLambda_{(i)}(\bx) = \text{diag}([\blambda_{(i)}(\bx)^\top \, 1]^\top)$, where a $1$ is appended at the end\footnote{Note that the homogeneous form adds a dummy neuron at each layer that is always active, since the last element of the hidden state is $1$ in homogeneous form, and thus always positive.} to match the dimensions of the homogeneous form in (\ref{Eq:NNHomogeneous}). Using the activation matrix, we can write the iteration defining the ReLU network from (\ref{Eq:NNHomogeneous}) as
\begin{align}
    \label{Eq:NNActivationDef}
    \begin{bmatrix}
    \bz_i\\
    1
    \end{bmatrix}
     = \bLambda_{(i)}(\bx)\bTheta_i\begin{bmatrix} \bz_{i-1}\\
    1
    \end{bmatrix}.
\end{align}
Finally, the map from input $\bx$ to the hidden layer preactivation $z_{ij}^\text{pre}$ can be written explicitly from the iteration (\ref{Eq:NNActivationDef}) as 
\begin{align}
\label{Eq:PreActExplicit}
    z_{ij}^\text{pre}(\bx) = \boldsymbol{\theta}_{ij}\Big(\prod_{l = 1}^{i-1}\bLambda_{(l)}(\bx)\bTheta_l\Big)\begin{bmatrix} \bx\\
    1
    \end{bmatrix},
\end{align}
where $\boldsymbol{\theta}_{ij}$ is the $j^{\text{th}}$ row of the matrix $\bTheta_i$. Similarly, the whole neural network function $\by = F_{\bTheta}(\bx)$ can be written explicitly as a map from $\bx$ to $\by$ as
\begin{align}
\label{Eq:NNExplicit}
    \by = \bTheta_L\Big(\prod_{i = 1}^{L-1}\bLambda_{(i)}(\bx)\bTheta_i\Big)\begin{bmatrix} \bx\\
    1
    \end{bmatrix}.
\end{align}
It is well known that in  (\ref{Eq:NNExplicit}), the output $\by$ is a continuous and PWA function of the input $\bx$ \cite{bengio,Hanin_2019,understanding}.  Next, we mathematically characterize PWA functions and give useful definitions.

\subsection{Polyhedra and PWA Functions}
\begin{definition}[(Convex) Polyhedron]
A convex polyhedron is a closed convex set $P \subset \mathbb{R}^n$ defined by a finite collection of $M$ halfspace constraints,
\[
P = \{\bx \mid \ba_i^\top \bx \le b_i, \quad i = 1, \ldots, M\},
\]
where $\ba_i,\bx \in \R^n$, $b_i \in \R$, and each $(\mathbf{a}_i, b_i)$ pair defines a halfspace constraint. Defining matrix $\mathbf{A} = [\mathbf{a}_1 \, \cdots \, \mathbf{a}_N]^\top$ and vector $\mathbf{b} = [b_1 \, \cdots \, b_N]^\top$, we can equivalently write
\begin{align*}
    P = \{\bx \mid \mathbf{Ax} \le \mathbf{b} \}. 
\end{align*}
We henceforth refer to convex polyhedra as just polyhedra. 
\end{definition}

The above representation of a polyhedron is known as the halfspace representation, or H-representation. 
% One can equivalently represent a polyhedron as the convex hull of a set of points and rays, known as the vertex representation, or V-representation, which we do not use in this paper. 
Note:
\begin{itemize}
\item A polyhedron can be bounded or unbounded.
\item A polyhedron can occupy a dimension less than the ambient dimension (if $(\mathbf{a}_i, b_i) = \alpha_{ij}(-\mathbf{a}_j, -b_j)$ for some pairs $(i,j)$ and some positive scalars $\alpha_{ij}$).
\item A polyhedron can be empty ($\nexists \bx$ such that $\mathbf{Ax} \le \mathbf{b}$).
\item Without loss of generality, the halfspace constraints for a polyhedron can be normalized so that $\|\ba_i\|\in \{0,1\}$ (by dividing the unnormalized parameters $\ba_i$ and $b_i$ by $\|\ba_i\|$), where $\|\cdot\|$ is the $\ell_2$ norm. $\|\ba\| = 0$ represents a degenerate constraint that is trivially satisfied if $b_i \ge 0$.
\item An H-representation of a polyhedron may have an arbitrary number of redundant constraints, which are constraints that are implied by other constraints. 
% Removing redundant constraints does not change the polyhedron.
\end{itemize}

\begin{definition}[Polyhedral Tessellation]\label{def:tessellation}
A polyhedral tessellation $\mathcal{P} = \{P_1, \ldots, P_N\}$ is a finite set of polyhedra that tessellate a set $\mathcal{X} \subset \mathbb{R}^n$. That is, $\cup_{i = 1}^NP_i = \mathcal{X}$ and $|P_i\cap P_j|_n = 0$, for all $i \not= j$, where $|\cdot|_n$ denotes the $n$-dimensional Euclidean volume of a set (the integral over the set with respect to the Lebesgue measure of dimension $n$). 
\end{definition}

Intuitively, the polyhedra in a tessellation together cover the set $\mathcal{X}$, and can only intersect with each other at shared faces, edges, vertices, etc. 
In the remainder of the paper when we refer to shared faces between neighboring polyhedra we mean the $n-1$ dimensional polyhedron given by the set intersection $P_i \cap P_j$.

\begin{definition}[Neighboring Polyhedra]
Two polyhedra, $P_i$ and $P_j$, are neighbors if their intersection has non-zero Euclidean volume in $n-1$ dimensions, that is, $|P_i \cap P_j|_{n-1} > 0$. 
\end{definition}

Intuitively, neighboring polyhedra are adjacent to one another in the tessellation, and share a common face. A polyhedral tessellation naturally induces a graph in which each node is a polyhedron and edges exist between neighboring polyhedra.

\begin{definition} [Piecewise Affine (PWA) Function] A PWA function is a function $F_{\text{PWA}}(\bx) : \mathbb{R}^{n} \rightarrow \mathbb{R}^{m}$ that is affine over each polyhdron in a polyhedral tessellation $\mathcal{P}$ of $\mathbb{R}^n$.  Specifically, $F_{\text{PWA}}(\bx)$ is defined as
\begin{align}
    F_{\text{PWA}}(\bx) = \mathbf{C}_k \bx + \mathbf{d}_k \ \ \ \forall \bx \in P_k \ \ \ \forall P_k \in \mathcal{P}, \label{eq:PWA}
\end{align}
where $\mathbf{C}_k \in \mathbb{R}^{m \times n}$, $\mathbf{d}_k \in \mathbb{R}^{m}$.
\end{definition}

Note that this requires a collection of tuples $\{(\mathbf{A}_1,\mathbf{b}_1, \mathbf{C}_1, \mathbf{d}_1), \ldots, (\mathbf{A}_N,\mathbf{b}_N, \mathbf{C}_N, \mathbf{d}_N)\}$, where $\mathbf{A}_k, \mathbf{b}_k$ define the polyhedron, and $\mathbf{C}_k, \mathbf{d}_k$ define the affine function over that polyhedron. 
% Here we consider only \emph{continuous} PWA functions, in which case the parameters are not independent, but must be such that each affine function continuously joins with the affine functions in each neighboring polyhedron.
We refer to the PWA representation in (\ref{eq:PWA}) as the \textit{explicit} PWA representation, as each affine map and polyhedron is written explicitly without specifying the relationship between them. 
There also exist other representations, such as the lattice representation \cite{lattice_orig,lattice_irred} and the, so called, canonical representation \cite{canonical, pwa_comp}. 
% Indeed, a ReLU network is also an example of a PWA representation in which the polyhedra and affine maps are implicitly contained in the neural network weights. 



\subsection{Dynamical Systems}
In this paper we are concerned with analyzing dynamical systems represented by ReLU networks.
Specifically, we consider the situations when the ReLU network represents the state transition function for a controlled (\ref{eq:controlled}) or autonomous (\ref{eq:uncontrolled}) dynamical system,
\begin{align}
    \bx_{t+1} &= F_{\bTheta}(\bx_t, \mathbf{u}_t) \label{eq:controlled} \\
    \bx_{t+1} &= F_{\bTheta}(\bx_t) \label{eq:uncontrolled}.
\end{align}
For notational convenience, in the remainder of the paper we drop the $\bTheta$ subscript from $F$.
We next define key concepts and state the overall problems that we seek to solve.

% Our goal in this paper is to transcribe the ReLU network function $\by = F_{\bTheta}(\bx)$ into its equivalent explicit PWA representation $\by = F_{\text{PWA}}(\bx)$, and then apply tools for PWA reachability analysis. 
\begin{definition}[$t$-Step Forward Reachable Set]
    Given an initial set of states $\mathcal{X}$ and a set of admissible control inputs $\mathcal{U}$, a $t$-step forward reachable set can be defined recursively,
    \begin{gather*}
    \mathcal{S}_0 = \mathcal{X} \\
        \mathcal{S}_{t} = \{\bx_{t} \mid \bx_{t} = F(\bx_{t-1}, \mathbf{u}_{t-1}),\\ \forall \bx_{t-1} \in \mathcal{S}_{t-1},\ \forall \mathbf{u}_{t-1} \in \mathcal{U}\}.
    \end{gather*}
\end{definition}

\begin{definition}[$t$-Step Backward Reachable Set]
    Given a final set of states $\mathcal{X}$ and a set of admissible control inputs $\mathcal{U}$, a $t$-step backward reachable set can be defined recursively,
    \begin{gather*}
    \mathcal{S}_0 = \mathcal{X} \\
        \mathcal{S}_{-t} = \{\bx_{-t} \mid \bx_{-t+1} = F(\bx_{-t}, \mathbf{u}_{-t}),\\ \forall \bx_{-t+1} \in \mathcal{S}_{-t+1},\ \forall \mathbf{u}_{-t} \in \mathcal{U}\}.
    \end{gather*}
\end{definition}

The ability to compute $t$-step forward and backward reachable sets will later enable us to compute control invariant sets and regions of attraction.
\begin{definition}[Control Invariant Set]
    Given a set of admissible control inputs $\mathcal{U}$, a control invariant set is a set of states for which there exists a sequence of control inputs such that the system state remains in the set for all time.
    If $\mathcal{S}$ is a control invariant set, then 
    \begin{gather*}
        \bx_0 \in \mathcal{S} \implies \exists \mathbf{u}_0, \ldots, \mathbf{u}_{t-1} \in \mathcal{U} \text{  s.t.  } \bx_t \in \mathcal{S} \quad \forall t \in \mathbb{N}
    \end{gather*}
    where $\bx_t = F(\bx_{t-1}, \mathbf{u}_{t-1})$ and $\mathbb{N}$ is the natural numbers.
\end{definition}

\begin{definition}[(Invariant) Region of Attraction]
    Given an autonomous dynamical system $\bx_{t+1} = F(\bx_t)$, an invariant region of attraction (ROA) is a set of states that asymptotically converges to an equilibrium state ($\bx_{\text{eq}}$) and for which sequences of states remain for all time. 
    If $\mathcal{S}$ is an ROA, then
    \begin{align*}
        \bx_0 \in \mathcal{S} &\implies \bx_t \in \mathcal{S} \quad \forall t \in \mathbb{N} \\
        \bx_0 \in \mathcal{S} &\implies \lim_{t \rightarrow \infty} \bx_t = \bx_{\text{eq}}
    \end{align*}
    where $\bx_t = F(\bx_{t-1})$.
\end{definition}

% \textcolor{black}{
% Due to the self-referential nature of the invariance properties for control invariant sets and ROAs, these sets tend to be more difficult to compute than forward or backward reachable sets.}

In this paper we seek to solve the following problems.
\begin{problem}
Given a ReLU network $\by = F(\bx)$ and a polyhedral input set $\mathcal{X} \subset \mathbb{R}^n$, compute the forward reachable set $\mathcal{Y}$. Similarly, given a polyhedral output set $\mathcal{Y}\subset \mathbb{R}^m$ compute the backward reachable set $\mathcal{X}$.
\end{problem}
\begin{problem}
Given a ReLU network that implements a discrete-time dynamical system $\bx_{t+1} = F(\bx_t)$ and state domain $\mathcal{X}$, identify the existence of stable fixed-point equilibria $\bx_{eq}$ and compute associated ROAs.
\end{problem}
\begin{problem}
Given a ReLU network that implements a discrete-time dynamical system $\bx_{t+1} = F(\bx_t, \mathbf{u}_t)$, a state domain $\mathcal{X}$, and set of admissible inputs $\mathcal{U}$, compute a control invariant set (if one exists).
\end{problem}
\begin{problem}
Given a ReLU network that implements a function $\by = F(\bx)$, determine whether the network is a homeomorphism (bijection with continuous inverse).
\end{problem}
In the next section we describe the RPM algorithm for transcribing the ReLU network function $\by = F_{\bTheta}(\bx)$ into its equivalent explicit PWA representation $\by = F_{\text{PWA}}(\bx)$. We then address the solution of the above problems in the subsequent sections.

%Our algorithm transforms a ReLU network into its explicit PWA representation by computing a polyhedral complex, $\mathcal{P} = \{P_1, \ldots, P_k, \ldots, P_N\}$, and associated affine maps $(\mathbf{C}_k, \mathbf{d}_k)$. The forward or backward reachable sets can then be computed separately for each polyhedron-affine map pair (using well-known methods from linear reachability theory), with the union giving the desired result for the full ReLU network. 

% Reachability of non-polyhedral sets can be performed by taking the union of individual polyhedral reachability analyses.